% Generated draft: requires mathematical and source review.
\documentclass[11pt]{article}
\usepackage[margin=25mm]{geometry}
\usepackage{fontspec}
\setmainfont{Latin Modern Roman}
\usepackage{amsmath,amssymb,mathtools,unicode-math}
\setmathfont{Latin Modern Math}
\usepackage{enumitem,needspace,xurl,hyperref}
\setlength{\parindent}{0pt}
\setlength{\parskip}{5pt}
\setlength{\emergencystretch}{3em}
\newcommand{\N}{\mathbb{N}}
\newcommand{\Z}{\mathbb{Z}}
\newcommand{\Q}{\mathbb{Q}}
\newcommand{\R}{\mathbb{R}}
\newcommand{\C}{\mathbb{C}}
\newcommand{\F}{\mathbb{F}}
\newcommand{\A}{\mathbb{A}}
\newcommand{\PP}{\mathbb{P}}
\newcommand{\E}{\mathbb{E}}
\newcommand{\eps}{\varepsilon}
\newcommand{\dd}{\,\mathrm{d}}
\begin{document}
\section*{UnsolvedMath ID: 30005548}
\textbf{Source problem code:} catalogue entry
\subsection{Definitions and mathematical statement}
For a positive integer \(d\), write \(\mathbb T^d=\R^d/\Z^d\), and let \(m_d\) denote its normalized Haar probability measure. The supplied excerpt asks whether a sequence denoted by \((\nu_n)\) converges to \(m_d\) in the weak-\(*\) topology and how fast or slowly this convergence occurs.

If the intended \(\nu_n\) are Borel probability measures on \(\mathbb T^d\), the convergence question has the precise meaning
\[
\nu_n\xrightarrow[n\to\infty]{\mathrm{weak}\text{-}*}m_d,
\qquad\text{that is,}\qquad
\lim_{n\to\infty}\int_{\mathbb T^d}f\,\dd\nu_n
=\int_{\mathbb T^d}f\,\dd m_d
\quad\text{for every }f\in C(\mathbb T^d;\R).
\]
This describes the convergence predicate only; it does not assert convergence for every sequence of measures. The excerpt does not define \(\nu_n\), the random walk, its acting group, its step distribution, its initial point or distribution, or the hypotheses under which convergence is sought. It also does not specify a distance between measures or a class and norm of test functions for a quantitative rate. Consequently, the intended quantified random-walk statement and its rate question cannot be recovered precisely from the supplied record.

The registry marks the statement as \texttt{unrecoverable}. Its supplied verification note says that the motivating question is answered immediately in the contribution by two stated semisimple random-walk theorems, and that no separate clean statement is promoted. Those theorem statements are not supplied here. The broader question concerning nondense orbits and general acting groups appears only in the supplied literature assessment and lacks sufficient hypotheses to serve as a separate mathematical statement.
\subsection{Short English statement}
The excerpt asks whether the measures of a torus random walk approach normalized Haar measure and at what rate. Weak-\(*\) convergence can be defined, but the measures and random-walk hypotheses are missing. The supplied record therefore does not determine a precise open problem; its partially solved label and accompanying assessment require human review.
\subsection{Sources}
\begin{itemize}
\item Local UnsolvedMath registry, \texttt{data/unsolvedmath/problems.json}, permanent ID \(30005548\), source problem code \texttt{catalogue entry}; title: ``Equidistribution Rates for Semisimple Torus Random Walks.'' The supplied record lists proposed year 2023 and no separate proposer.
\item Supplied contribution attribution: Weikun He, joint work with Nicolas de Saxc\'e, ``Semisimple linear random walks on the torus, II,'' p.~1431, in \emph{Incidence Problems in Harmonic Analysis, Geometric Measure Theory, and Ergodic Theory} (2023), Oberwolfach Reports, pp.~1397--1452. Supplied report DOI: \url{https://doi.org/10.4171/owr/2023/25}. The background describes Oberwolfach Reports as documenting mathematical workshops and their open-problem discussions.
\item Supplied literature source, ``OWR 2023 25'': \url{https://publications.mfo.de/bitstream/handle/mfo/4055/OWR_2023_25.pdf?sequence=1}.
\item Supplied literature source, ``He de Saxc\'e'': \url{https://doi.org/10.1017/etds.2025.3}.
\item Literature checked at: \texttt{2026-08-22}. This is supplied dated curation derived from the OWR record, not an independently verified assessment. The supplied status is ``partially solved.'' The supplied assessment says that He and de Saxc\'e prove quantitative equidistribution for semisimple Zariski closure, while a broader random-walk question involving nondense orbits and general acting groups is not fully resolved. The record expressly calls for independent verification of time-sensitive status and cited claims.
\item The supplied statement-verification note says that two semisimple random-walk theorems in the contribution answer the motivating question immediately, and that the original record remains publicly visible for provenance without a promoted clean statement. No source was independently consulted for this draft.
\end{itemize}
\end{document}
